%=============================================================================
%  CHAPTER 7 — TESTING, VALIDATION, AND DEPLOYMENT
%=============================================================================
\chapter{Testing, Validation, and Deployment}

\section*{Chapter Introduction}

This chapter presents the testing and validation process of the ForsaDrive system. It aims to demonstrate that the implemented solution meets the functional and non-functional requirements defined earlier. The chapter describes the testing strategy, presents representative functional test cases, evaluates system performance and compatibility, and outlines the deployment approach. Finally, it summarises the results obtained and assesses the overall reliability of the system.

%─────────────────────────────────────────────────────────────────────────────
\section{Testing Strategy}

To ensure system correctness and reliability, a structured testing strategy was adopted. The testing process combined both manual and tool-assisted validation methods.

The following categories of testing were performed:

\begin{itemize}[noitemsep]
    \item \textbf{Functional Testing:} Verification of system features against requirements
    \item \textbf{Integration Testing:} Validation of interactions between backend, web, and mobile components
    \item \textbf{API Testing:} Verification of REST endpoints using Postman
    \item \textbf{User Interface Testing:} Evaluation of usability and responsiveness
    \item \textbf{Compatibility Testing:} Testing across different devices and browsers
\end{itemize}

Testing was conducted iteratively throughout development to detect and correct issues early.

%─────────────────────────────────────────────────────────────────────────────
\section{Functional Testing}

Functional testing ensures that the system behaves as expected from the user's perspective.

\begin{table}[H]
\centering
\caption{Sample Functional Test Cases}
\label{tab:functional-tests}
\begin{tabularx}{\textwidth}{c l X X X}
\toprule
\textbf{ID} & \textbf{Feature} & \textbf{Test Scenario} & \textbf{Expected Result} & \textbf{Result} \\
\midrule
TC01 & Registration & Valid user data entered & Account created successfully & Passed \\
TC02 & Login & Correct credentials & User authenticated & Passed \\
TC03 & Publish Ride & Valid ride data & Ride visible in system & Passed \\
TC04 & Book Ride & Available seats + sufficient balance & Booking confirmed & Passed \\
TC05 & Messaging & Send message & Message delivered & Passed \\
TC06 & Wallet & Booking deduction & Balance updated correctly & Passed \\
\bottomrule
\end{tabularx}
\end{table}

These tests confirm that core functionalities are correctly implemented.

%─────────────────────────────────────────────────────────────────────────────
\section{Requirements Validation}

Each requirement defined in Chapter 3 was validated through implementation and testing.

\begin{itemize}[noitemsep]
    \item Ride management features were validated through publish/search/book workflows
    \item Authentication was validated through login and session control
    \item Wallet operations were validated through booking transactions
    \item Messaging and rating systems were verified through user interaction scenarios
\end{itemize}

This ensures traceability between requirements and implementation.

%─────────────────────────────────────────────────────────────────────────────
\section{API Testing}

The backend APIs were tested using Postman to ensure correct request handling and response structure.

\begin{itemize}[noitemsep]
    \item Validation of request parameters
    \item Verification of response codes (200, 400, etc.)
    \item Testing of authentication tokens
\end{itemize}

Example API response structure:

\begin{verbatim}
{
  "status": "success",
  "data": {...}
}
\end{verbatim}

This confirms that the backend provides consistent and predictable responses.

%─────────────────────────────────────────────────────────────────────────────
\section{Compatibility and Performance Testing}

\subsection{Compatibility Testing}

The system was tested on multiple platforms:

\begin{itemize}[noitemsep]
    \item Web browsers: Chrome, Firefox
    \item Mobile devices: Android smartphones
\end{itemize}

Results showed consistent behaviour across platforms.

\subsection{Performance Testing}

Basic performance evaluation was conducted:

\begin{itemize}[noitemsep]
    \item Response time remained within acceptable limits
    \item No major delays observed under normal usage
\end{itemize}

Due to project constraints, large-scale load testing was not performed.

%─────────────────────────────────────────────────────────────────────────────
\section{Security Testing}

Security testing focused on protecting the system against common vulnerabilities:

\begin{itemize}[noitemsep]
    \item Input validation to prevent injection attacks
    \item Authentication checks for protected endpoints
    \item Role-based access control verification
\end{itemize}

These measures contribute to system reliability and data protection.

%─────────────────────────────────────────────────────────────────────────────
\section{Deployment Process}

The system was deployed in a local development environment.

\subsection{Backend Deployment}

\begin{itemize}[noitemsep]
    \item PHP server setup
    \item Database initialisation
\end{itemize}

\subsection{Web Application}

\begin{itemize}[noitemsep]
    \item Accessible via local server
\end{itemize}

\subsection{Mobile Application}

\begin{itemize}[noitemsep]
    \item Deployed on Android devices via APK
\end{itemize}

Although production deployment was not implemented, the system is designed to be easily migrated to a hosted environment.

%─────────────────────────────────────────────────────────────────────────────
\section{Results and Evaluation}

The testing process demonstrated that:

\begin{itemize}[noitemsep]
    \item All core functionalities are operational
    \item The system is stable under normal usage
    \item User interactions are smooth and intuitive
    \item Web and mobile platforms behave consistently
\end{itemize}

These results validate the effectiveness of the implemented solution.

%─────────────────────────────────────────────────────────────────────────────
\section*{Chapter Conclusion}

This chapter presented the testing and validation of the ForsaDrive system, confirming that it meets its functional requirements and performs reliably under expected conditions. The deployment process was also outlined, demonstrating the system's readiness for future production deployment.

The next chapter concludes the report by summarising the project achievements and presenting future perspectives.